\section*{Capítulo \ref{ch:g5:raw-materials-and-recycling} --- De las materias primas al reciclaje}

\begin{solution}{exo:g5:raw-materials-and-recycling:1}
Naturales: la madera, la lana, la piedra. Elaborados: el vidrio, el
acero, el plástico.
\end{solution}

\begin{solution}{exo:g5:raw-materials-and-recycling:2}
El papel: madera (árboles). El vidrio: arena (con sosa y caliza). El
algodón: la planta del algodón. El plástico: petróleo.
\end{solution}

\begin{solution}{exo:g5:raw-materials-and-recycling:3}
Una roca de la que se extrae un metal. El aluminio viene de la bauxita.
\end{solution}

\begin{solution}{exo:g5:raw-materials-and-recycling:4}
El tarro, en el contenedor de vidrio; el periódico, con el papel y el
cartón; la lata, con el metal; la botella, con las botellas de plástico.
\end{solution}

\begin{solution}{exo:g5:raw-materials-and-recycling:5}
\omterm{def:g5:raw-materials-and-recycling:raw-material}{Materia prima} nueva: el material \omterm{def:g5:raw-materials-and-recycling:recycling}{reciclado} se usa en lugar de la arena,
la \omterm{def:g5:raw-materials-and-recycling:raw-material}{mena} o el petróleo tomados de la naturaleza.
\end{solution}

\begin{solution}{exo:g5:raw-materials-and-recycling:6}
Renovables: la madera (si se replanta), la lana, el algodón. No
renovables: el petróleo, la \omterm{def:g5:raw-materials-and-recycling:raw-material}{mena} de hierro.
\end{solution}

\begin{solution}{exo:g5:raw-materials-and-recycling:7}
Fabricar vidrio a partir de arena es un \omterm{def:g5:irreversible-changes:chemical-change}{cambio químico}: se fabrica algo
nuevo, y no hay un camino de vuelta sencillo.
\end{solution}

\begin{solution}{exo:g5:raw-materials-and-recycling:8}
Por ejemplo: reducir, comprar fruta sin envoltorio de plástico;
reutilizar, rellenar un tarro de vidrio; reciclar, tirar una lata de
refresco al contenedor del metal.
\end{solution}

\begin{solution}{exo:g5:raw-materials-and-recycling:9}
Cada material se recicla a su manera: trozos de otro material
estropearían el vidrio fundido.
\end{solution}

\begin{solution}{exo:g5:raw-materials-and-recycling:10}
Una veinteava parte de 20 unidades: $20 \div 20 = 1$ unidad de energía,
más o menos.
\end{solution}

\begin{solution}{exo:g5:raw-materials-and-recycling:11}
Reciclar sigue necesitando energía, máquinas y transporte; un objeto
que nunca se fabrica ahorra su \omterm{def:g5:raw-materials-and-recycling:raw-material}{materia prima} y toda esa energía.
\end{solution}
